\section{Conclusion}
\label{sec:conclusion}

This paper introduced, mathematically formalized, and empirically validated the \textbf{Cloud-Native Observability Maturity Model (CN-OMM)}. Developed in strict adherence to Design Science Research guidelines (Becker et al.'s R1--R8 \cite{becker_2009}, De Bruin et al.'s lifecycle phases \cite{debruin_2005}, and Lasrado et al.'s configurational theory \cite{lasrado_2015}), the CN-OMM resolves three foundational weaknesses in prior industrial and academic frameworks: vendor lock-in bias, blindness to the instrumentation probe effect, and the monomodal additivity fallacy.

By structuring observability across a $5\text{-Level} \times 5\text{-Dimension}$ stage-gate matrix ($D_1$: Telemetry Acquisition \& Depth, $D_2$: Semantic Context \& Correlation, $D_3$: Diagnostic Intelligence, $D_4$: Resilience Verification \& SLOs, and $D_5$: Instrumentation Economics) and operationalizing 15 automated Quantitative Capability Indicators (QCIs), the CN-OMM replaces subjective questionnaire assessments with deterministic, programmatic cluster auditing. Our Type-3 empirical validation over $N=30$ paired Chaos Engineering trials on Kubernetes proved that advancing to monomodal distributed tracing (Level 3) without fulfilling the configurational stage-gate of semantic cross-pillar correlation ($D_2$) collapses root-cause accuracy to $0.00\%$ due to \textit{Cascading Span Bias}. Conversely, achieving balanced Level 4 maturity (multimodal correlation with exclusive self-time decomposition $T_{\text{self}}$) elevates diagnostic accuracy to $83.33\%$ ($p = 7.74 \times 10^{-7}$) and reduces median diagnostic latency by $94.24\%$ ($300.00\,\text{s} \rightarrow 17.27\,\text{s}$, $p = 3.06 \times 10^{-5}$) while bounding steady-state CPU overhead at $9.51\% \pm 2.14\%$ ($p = 6.52 \times 10^{-9}$).

\textbf{Open Science and Reproducibility:} To support full scientific auditability (Requirement R8), our complete replication package---including the LaTeX source, the automated CN-OMM cluster assessment engine (\texttt{scripts/evaluate\_cn\_omm.py}), the capability audit datasets, and the paired Chaos Engineering trial logs---is publicly available in our repository. Future work will extend the PANOPTES reference implementation from Level 4 to Level 5 by integrating neuro-symbolic Graph Neural Networks and autonomous Kubernetes remediation controllers governed by dynamic FinOps/GreenOps signal weighting.
